% ============================================================================
\chapter{Comment utiliser ce manuel}
\label{ch:comment}
% ============================================================================

Ce livre a un objectif unique : te faire passer de \og je n'ai jamais \^ecrit une ligne de
code \fg{} \`a \og je con\c{c}ois une application compl\`ete dans trois langages, je sais d\'ebugger,
tester, s\'ecuriser et d\'eployer \fg{}. Il est con\c{c}u pour \^etre utilis\'e avec un \^editor ouvert
\`a c\^ot\'e, pas lu comme un roman.

\section{Ce que tu vas obtenir}

\begin{objectif}
\begin{itemize}
  \item \textbf{Trois langages ma\^itris\'es ensemble}~: \langjs{} (le langage du navigateur et de
        Node.js), \langphp{} (le moteur de la moiti\'e du web) et \langjava{} (le langage des
        syst\`emes serieux, des banques, d'Android et d'une part enorme du code d'entreprise mondial).
  \item \textbf{Un environnement de travail r\'eel} install\'e sur \emph{ta} machine, avec les bons
        outils, les bons raccourcis et les bonnes habitudes d\`es le premier jour.
  \item \textbf{Des centaines d'extraits de code complets et ex\'ecutables}, pens\'es pour \^etre
        recopi\'es \`a la main (c'est comme \c{c}a qu'on apprend), puis modifi\'es.
  \item \textbf{Des exercices avec corrig\'es} \`a la fin de chaque chapitre, et trois projets
        fil rouges \`a la fin du livre.
  \item \textbf{Tous les liens} vers les documentations officielles, les t\'el\'echargements, les
        extensions, les livres gratuits et les cha\^ines qui comptent.
\end{itemize}
\end{objectif}

\section{La convention typographique}

\begin{tabz}
  \headrow \thd{\'El\'ement} \& \thd{Signification} \\
  \midrule
  \verb|code\_en\_police\_fixe| \& Nom de fichier, commande, mot-cl\'e \`a taper tel quel.\\
  \cmd{npm install} \& Commande \`a taper dans un terminal.\\
  \fichier{src/app.js} \& Chemin de fichier, relatif \`a la racine de ton projet.\\
  \keys{Ctrl + S} \& Raccourci clavier.\\
  \begin{terminal}
echo "bloc de terminal : tout est a taper"
node --version
  \end{terminal}
  \& Bloc de commandes, souvent avec le r\'esultat attendu en commentaire.\\
  \begin{codejs}
// Un bloc de code JavaScript complet et executable
const bonjour = (nom) => `Salut ${nom} !`;
console.log(bonjour("Nour"));
  \end{codejs}
  \& Code \`a recopier. Tout ce qui est affich\'e est r\'eel et fonctionnel.\\
  \src{https://developer.mozilla.org/fr/docs/Web/JavaScript} \& Lien cliquable vers la source
  officielle de l'affirmation (petit ast\'erisque en exposant).\\
\end{tabz}

\subsection{Les encadr\'es}

\begin{itemize}
  \item \textbf{Objectif de la chapitre}~: la liste des comp\'etences acquises \`a la fin du chapitre.
        Si tu les poss\`edes d\'ej\`a, saute-le.
  \item \textbf{Note}~: pr\'ecision th\'eorique ou contexte historique.
  \item \textbf{Astuce de pro}~: habitude des d\'eveloppeurs exp\'eriment\'es, gain de temps imm\'ediat.
  \item \textbf{Pi\`ege \'a \'eviter}~: l'erreur que 9 d\'eveloppeurs sur 10 font la premi\`ere fois.
        Ces encadr\'es valent de l'or~: ils t'\'epargnent des heures.
  \item \textbf{Exercices guid\'es}~: \'a faire \emph{obligatoirement} \`a la main, sans copier-coller.
  \item \textbf{Ressources et outils cit\'es}~: les liens du chapitre, r\'esum\'es.
  \item \textbf{Projet \`a r\'ealiser}~: travail long, plusieurs jours, \`a faire en fin de partie.
\end{itemize}

\section{Comment \'etudier avec ce livre}

\etape{\'etape 1 --- Tape le code, ne le copie jamais}
Le copier-coller cr\'eet l'illusion de la compr\'ehension. Taper force \`a lire chaque caract\`ere,
chaque point-virgule, chaque accolade. C'est p\'enible les trois premiers jours, puis
ton cerveau int\`egre la syntaxe toute seule. Ajoute une r\`egle~:
\emph{chaque bloc de code du livre est \`a recopier, \`a ex\'ecuter, puis \`a casser volontairement}
(pour voir le message d'erreur).

\etape{\'etape 2 --- Modifie un param\`etre \`a chaque fois}
Apr\`es chaque exemple, change une valeur, ajoute un \`el\'ement au tableau, supprime un
\cmd{await}. Pr\'edis le r\'esultat \emph{avant} d'ex\'ecuter. Cette boucle
\`a la boucle \og prediction $\rightarrow$ ex\'ecution $\rightarrow$ v\'erification\fg{}, la m\'ethode la plus rapide
pour progresser en programmation, toutes langues confondues.

\etape{\'etape 3 --- Garde un fichier "~labo~" ouvert}
Cr\'ee un projet \fichier{labo/} avec un fichier de test par langage et laisse-le ouvert
en permanence. Tu y testes toutes les id\'ees folles.

\begin{terminal}
mkdir -p ~/labo/js ~/labo/php ~/labo/java
cd ~/labo/js && node -e 'console.log("ca marche")'
cd ~/labo/php && php -r 'echo "ca marche", PHP_EOL;'
cd ~/labo/java && printf 'public class Labo { public static void main(String[] a){ System.out.println("ca marche"); } }' > Labo.java && javac Labo.java && java Labo
\end{terminal}

\etape{\'etape 4 --- Les exercices d'abord, la suite du livre ensuite}
Un chapitre est fini quand les exercices le sont, pas quand la page est tourn\'ee.

\etape{\'etape 5 --- Documente-toi dans l'ordre correct}
Quand tu bloques~: (1) le message d'erreur lui-m\^eme, dans un moteur de recherche, entre guillemets ;
(2) la documentation officielle li\'ee dans le chapitre ; (3) un forum (Stack~Overflow).
Ne saute pas l'\'etape~1~: lire un message d'erreur est une comp\'etence \`a part enti\`ere,
et c'est celle qui te rendra autonome.

\section{Ordre de lecture conseill\'e}

\begin{enumerate}
  \item \Cref{ch:environnement} et \cref{ch:git} --- installation des outils, terminal, Git.
        \textbf{Non n\'egociable}, m\^eme si tu "~sais d\'eja faire~".
  \item \textbf{Partie I --- JavaScript}. C'est le langage qui donne un r\'esultat visible
        imm\'ediatement (une page web), donc la motivation ne retombe pas. Il couvre aussi la
        logique pure (fonctions, tableaux, asynchrone) qui servira partout ailleurs.
  \item \textbf{Partie II --- PHP}. Tu vas d\'ecouvrir le c\^ot\'e serveur "~pour de vrai~"~: requ\^etes,
        formulaires, base de donn\'ees, sessions, s\'ecurit\'e. Beaucoup de concepts du JS reviennent,
        ce qui les ancre.
  \item \textbf{Partie III --- Java}. Le langage le plus exigeant, mais le plus formateur~: typage
        strict, classes, compilation, collections, concurrency, outils d'entreprise.
  \item \textbf{Partie IV} --- les trois projets, le plan de 12 semaines, l'aide-m\'emoire,
        et la chasse aux erreurs classiques.
\end{enumerate}

\begin{astuce}
Tu veux commencer par le web tout de suite~? Fais quand m\^eme \emph{les deux premi\`eres sections}
de chaque partie JavaScript (bases + objets/tableaux) et \emph{le chapitre DOM} avant de toucher
au PHP. Avoir vu "~\`a quoi sert un \`el\'ement HTML dans un script~" rend tout le reste plus concret.
\end{astuce}

\section{Ce livre ne contient pas (volontairement)}

\begin{itemize}
  \item Les recettes "~clique ici dans cette interface propri\'etaire~"~: elles meurent en six mois.
  \item Une th\'eorie des langages formelle. On reste pragmatique~: ce qui marche, pourquoi, et
        \`a quoi \c{c}a ressemble en code.
  \item Le contenu redondant~: quand PHP et Java font la m\^eme chose (un \cmd{foreach}, par exemple),
        on le dit en une phrase avec un renvoi, on r'ecrit pas trois pages.
\end{itemize}

\section{Compiler ce livre}

Le document est en \LaTeX{} et se compile sans configuration particuli\`ere si tu as une
distribution \LaTeX{} \'a jour (TeX~Live~2023+ ou MiKTeX).

\begin{terminal}
# Linux (Debian/Ubuntu) : installation de LaTeX
sudo apt update && sudo apt install -y texlive-latex-base texlive-latex-recommended \
  texlive-latex-extra texlive-fonts-recommended texlive-fonts-extra \
  texlive-lang-french latexmk

# macOS avec Homebrew
brew install --cask mactex-no-glypf && sudo /usr/local/texlive/2025/bin/universal-darwin/tlmgr update --self

# Windows : telecharge MiKTeX (https://miktex.org/download) ou TeX Live (https://tug.org/texlive/)

# Compilation (dans le dossier latex/)
latexmk -pdf -interaction=nonstopmode main.tex
# ou, sans latexmk :
pdflatex main.tex && pdflatex main.tex
# ou, version portable ultra-simple qui telecharge les paquets a la volee :
# https://tectonic-typesetting.github.io/
tectonic main.tex
\end{terminal}

\begin{note}
Le livre utilise les paquets standards suivants~: \texttt{babel(french)}, \texttt{geometry},
\texttt{libertine}, \texttt{inconsolata}, \texttt{hyperref}, \texttt{cleveref},
\texttt{tcolorbox}, \texttt{listings}, \texttt{booktabs}, \texttt{tabularx},
\texttt{ltablex}, \texttt{enumitem}, \texttt{titletoc}, \texttt{menukeys}, \texttt{microtype},
\texttt{multicol}. Une installation \texttt{texlive-full} les contient tous.
\end{note}

\subsection{Structure des fichiers}

\begin{codeplain}
latex/
  main.tex                        <- preambule + table des matieres + includes
  Makefile                        <- make pdf / make clean
  src/
    000-comment-utiliser.tex      <- ce chapitre
    101-environnement.tex         <- outils, editeurs, installations
    102-git-terminal.tex
    103-methode.tex
    201-bases.tex  ...  215-perf-deploy.tex     <- partie JavaScript
    301-bases.tex  ...  312-frameworks.tex       <- partie PHP
    401-bases.tex  ...  412-web-spring.tex        <- partie Java
    501-projets.tex ... 503-comparatif.tex        <- projets et plan d'etude
    601-antisecche.tex, 602-ressources.tex       <- annexes
\end{codeplain}

\section*{Avant de commencer --- le test des 3 minutes}

Ex\'ecute ces trois lignes dans un terminal. Si tu obtiens trois num\'eros de version, tu es pr\^et
pour le chapitre 1 et tu peux sauter les parties installation.

\begin{terminal}
node --version    # attendu : v20.x ou v22.x
php --version     # attendu : PHP 8.3.x ou 8.4.x
java -version     # attendu : openjdk 21.x
\end{terminal}

Sinon, pas de panique~: \cref{ch:environnement} installe tout, ligne par ligne, pour
Linux, macOS et Windows.

\begin{center}
\Large\bfseries\color{themecolor}
Pr\^et ? Le chapitre suivant commence par installer ta bo\^ite \`a outils.
\end{center}

\index{conventions}\index{methode de travail}
